\section{The finite-memory frontier beyond reported resetting}
\label{sec:r79-quantum-frontier}

Reported resetting permits the scalar and memory-preserving reduction
of Theorem~\ref{thm:r78-palette-normal-form}.  The next result does not
assume resetting.  It permits arbitrary finite-dimensional failure
maps and an unknown model which remains fixed during the experiment.
The appropriate finite certificate retains the counted controller
labels and their reachable quantum supports.  In particular, the
quantum state need not be determined by the latest report.

The support construction and its potential equations are inherited
from Lemma~\ref{lem:reachable-support},
Theorems~\ref{thm:memory-certificate} and~\ref{thm:face}; finite-model
algebraic synthesis is Corollary~\ref{cor:memory-algebraic}.
The primitive fallback is already supplied by
Theorem~\ref{thm:primitive-fallback} and
Lemma~\ref{lem:compact-fallback}.  We give their positive-cost form
in full to make the new stability argument self-contained.  The
extension here permits a compact semialgebraic model family with
unrestricted legal feedback.  Reachability and cost observables
are established ingredients
\cite{YuYing12,YingFengYuYing13,LiuZhouBartheYing25,AvanziniEtAl22},
as are classical bounded-controller and robust synthesis
\cite{AmatoBernsteinZilberstein10,JungesEtAl18,CubuktepeEtAl21,
GaleslootEtAl25}.  The next section supplies the new comparison
that preserves the exact memory budget.

\subsection{Primitive instruments and counted controllers}

Fix a Hilbert space $\mathcal H$ of dimension $d\ge1$, a nonempty
finite command set $\mathcal A$, and a finite nonterminal report set
$\mathcal Y$.  Let $\Theta$ be a nonempty compact parameter space.
For every $\theta\in\Theta$ and $a\in\mathcal A$, the nonterminal
branches are completely positive maps
$\Phi^\theta_{a,y}:\mathcal B(\mathcal H)\to
\mathcal B(\mathcal H)$ satisfying
\[
 B^\theta_a=I-\sum_{y\in\mathcal Y}
                    (\Phi^\theta_{a,y})^*(I)\succeq0.
\]
The effect $B^\theta_a$ is the probability of success; success
requires immediate stopping and no other stopping is permitted.
The mean call-cost operator is $G^\theta_a$, with common constants
\begin{equation}\label{eq:r79-positive-costs}
 0<c_*I\preccurlyeq G^\theta_a\preccurlyeq c^*I<\infty.
\end{equation}
The maps and costs depend continuously on $\theta$.  All commands
are legal, and the only observations are the sampled command, the
declared report, and success.  The realized cost is not an additional
observation.  The physical implementation of these maps and costs is
part of the instrument specification, as in
Section~\ref{sec:reset-control}.

The adversary chooses one pair $(\theta,\rho)$, consisting of a
model and an initial density matrix.  This pair is not observed by
the controller.  The same $\theta$ governs every subsequent call.
Thus the model is not reselected after a report, and the observed
history can contain information about it.

An autonomous controller with at most $M\ge1$ live labels is specified
by command rows and update rows
\begin{equation}\label{eq:r79-controller-parameters}
 p^k\in\Delta(\mathcal A),\qquad
 u^{k,a,y}\in\Delta([M]),\qquad 1\le k\le M.
\end{equation}
Its prescribed initial label is $1$.  At label $k$ it samples $a$
from $p^k$; after nonterminal report $y$ it writes label $\ell$ with
probability $u^{k,a,y}_\ell$.  These are fresh random choices, not an
uncounted persistent seed.  Unused labels and duplicated rows are
allowed.  Denote this compact product of simplexes by $\mathcal C_M$.

For a controller $\gamma=(p,u)$, work in the direct-sum matrix algebra
\[
 \mathfrak A_M=\bigoplus_{k=1}^M\mathcal B(\mathcal H),
 \qquad
 \langle W,X\rangle=\sum_k\tr(W_kX_k).
\]
The nonterminal evolution and call-cost operator are
\begin{align}
 (\mathsf T^\theta_\gamma X)_\ell
   &=\sum_{k,a,y}p^k_a u^{k,a,y}_\ell
                         \Phi^\theta_{a,y}(X_k),
       \label{eq:r79-product-map}\\
 K^\theta_{\gamma,k}&=\sum_a p^k_aG^\theta_a.
       \label{eq:r79-product-cost}
\end{align}
The map $\mathsf T^\theta_\gamma$ is completely positive and
trace-nonincreasing.  Let $\iota_1(\rho)=(\rho,0,\ldots,0)$.
The total expected cost and its robust value are
\begin{align}
 J_\gamma(\theta,\rho)
   &=\sum_{n=0}^\infty
       \langle K^\theta_\gamma,
          (\mathsf T^\theta_\gamma)^n\iota_1(\rho)\rangle,
          \label{eq:r79-total-cost}\\
 J(\gamma)&=\sup_{\theta\in\Theta,\ \rho\succeq0,\ \tr\rho=1}
                       J_\gamma(\theta,\rho),
 \qquad V_M=\inf_{\gamma\in\mathcal C_M}J(\gamma).
          \label{eq:r79-general-frontier}
\end{align}
All summands are nonnegative.  Consequently this definition is
unambiguous even for an improper controller.  Finite cost implies
finite expected duration by \eqref{eq:r79-positive-costs}.

\subsection{A primitive feasibility criterion}

For a fixed model, choose any Kraus representations
$\Phi^\theta_{a,y}(X)=\sum_\mu
 A^\theta_{a,y,\mu}X(A^\theta_{a,y,\mu})^*$.
A nonzero subspace $S\subseteq\mathcal H$ is \emph{closed and dark}
if
\begin{equation}\label{eq:r79-dark-subspace}
 A^\theta_{a,y,\mu}S\subseteq S\quad\hbox{for all }a,y,\mu,
 \qquad B^\theta_a|_S=0\quad\hbox{for all }a.
\end{equation}
The condition is representation-independent: it says that every
command has zero success probability on $S$ and all its failure
branches preserve $S$.

\begin{theorem}[Feasibility from arbitrary initial quantum states]
\label{thm:r79-dark-feasibility}
The following conditions are equivalent.
\begin{enumerate}
\item There is one visible-history policy with finite uniform
expected cost for every $(\theta,\rho)$.
\item No model $\theta\in\Theta$ has a nonzero closed dark subspace.
\item Every fixed full-support command distribution $\bar p$,
used independently on every call with one live label, has finite
uniform expected cost for every $(\theta,\rho)$.
\end{enumerate}
In particular, feasibility for unrestricted history and feasibility
at every finite live-label budget coincide.  For such a $\bar p$ set
\[
 \mathsf F^\theta=\sum_{a,y}\bar p_a\Phi^\theta_{a,y},
 \qquad \bar G^\theta=\sum_a\bar p_aG^\theta_a.
\]
Then $r(\mathsf F^\theta)<1$ for every $\theta$, and an explicit
uniform cost bound is
\begin{equation}\label{eq:r79-fallback-cap}
 C_0=\max_{\theta\in\Theta}
       \lambda_{\max}\!\left[
          (I-(\mathsf F^\theta)^*)^{-1}\bar G^\theta\right]<\infty.
\end{equation}
\end{theorem}
\begin{proof}
If \eqref{eq:r79-dark-subspace} holds, an initial density matrix
supported on $S$ remains there after every failure and has zero
success probability at every call, regardless of the policy.  Its
cost is infinite by the positive floor.  This proves the necessity
of the second condition.

For sufficiency fix $\theta$ and write $\mathsf F=\mathsf F^\theta$.
It is positive and trace-nonincreasing, so its spectral radius is
at most one.  One direct justification is that, for Hermitian $X$,
$-\|X\|I\preccurlyeq X\preccurlyeq\|X\|I$ and positivity bounds
$\mathsf F^n(X)$ by $\|X\|\mathsf F^n(I)$ in order; the latter
has trace at most $d$.  Decomposing a general matrix into its real
and imaginary Hermitian parts shows that the powers are bounded.

Suppose its spectral radius is one.  The numbers
$\tr(\mathsf F^n(I))$ decrease.  Their limit is positive: otherwise
the same order bounds would make $\mathsf F^n$ tend to zero on
every matrix, contradicting the existence of an eigenvalue of
modulus one.  The positive Ces\`aro averages
\[
 Z_N=\frac1N\sum_{n=0}^{N-1}\mathsf F^n(I)
\]
have bounded trace and a subsequence converging to $Z\succeq0$ with
$\tr Z>0$.  Since
$\mathsf F(Z_N)-Z_N=(\mathsf F^N(I)-I)/N$, one has
$\mathsf F(Z)=Z$.

Put $S=\operatorname{ran}Z$.  Taking traces in this equality gives
$\sum_a\bar p_a\tr(B^\theta_aZ)=0$.  Every term is nonnegative
and every $\bar p_a$ is positive, so $B^\theta_a|_S=0$ for every
$a$.  The identity $\mathsf F(Z)=Z$ is a sum of positive Kraus
images.  Pairing it with a vector in $\ker Z$ shows, term by term,
that every Kraus image has range in $S$; equivalently
$A^\theta_{a,y,\mu}S\subseteq S$.  This constructs a nonzero
closed dark subspace.  Therefore its absence forces
$r(\mathsf F^\theta)<1$.

The geometric operator series converges and gives the cost operator
in \eqref{eq:r79-fallback-cap}.  The primitive maps and costs are
continuous, and the inverse is continuous wherever it exists.
Compactness of $\Theta$ therefore makes the displayed maximum
finite.  This single row works for all models and all initial
density matrices, proving the third condition and hence the first.
\end{proof}

The all-initial-state quantifier is material to the necessity:
a particular supplied initial state might never reach a closed dark
subspace.  The theorem also uses the fixed finite alphabet and the
availability of every command at every call.

\subsection{Reachable corners and exact cost witnesses}

For fixed $\theta,\gamma$, abbreviate $\mathsf T=\mathsf
T^\theta_\gamma$ and $K=K^\theta_\gamma$.  A tuple of orthogonal
projections $P=(P_1,\ldots,P_M)$ defines the matrix corner
$P\mathfrak A_MP=\bigoplus_kP_k\mathcal B(\mathcal H)P_k$.
Call it invariant if
\begin{equation}\label{eq:r79-corner-invariance}
 P_1=I,\qquad P\mathsf T(P)P=\mathsf T(P).
\end{equation}
Multiplication in this formula is blockwise.  By positivity, this
identity implies that $\mathsf T$ maps the entire corner into
itself.  Indeed every positive element of the corner is bounded
above by a scalar multiple of $P$, and general elements follow by
linear decomposition.

\begin{lemma}[Finite reachable-support construction]
\label{lem:r79-reachable-corner}
The least invariant corner containing the initial label is obtained
by taking the blockwise support of
\begin{equation}\label{eq:r79-reachable-sum}
 R=\sum_{n=0}^{Md-1}\mathsf T^n\iota_1(I).
\end{equation}
In particular it is determined after finitely many linear operations
and support computations, without an infinite-history enumeration.
If $J_\gamma(\theta,\rho)$ is finite for every initial density matrix,
the restriction of $\mathsf T$ to this corner is transient and its
cost operator is finite on every block of the corner.
\end{lemma}
\begin{proof}
Let $R_n=\sum_{j=0}^n\mathsf T^j\iota_1(I)$ and let $P^{(n)}$
be its blockwise support.  These supports increase, beginning with
total rank $d$ and having total rank at most $Md$.  If no rank is
added at a step, then $\mathsf T(R_n)$ is supported on $P^{(n)}$.
The operator $R_n$ is positive definite on this support, so
$P^{(n)}$ is bounded by a finite multiple of $R_n$.  Positivity now
implies \eqref{eq:r79-corner-invariance}, and all later supports
coincide.  There can be at most $Md-d$ strict enlargements, so
\eqref{eq:r79-reachable-sum} has already reached the least invariant
corner.  Every invariant corner with first block $I$ contains all
the summands, proving minimality.

For the last assertion, finite cost for every initial density matrix
implies finite cost for $\iota_1(I)$.  For each fixed $n$, the total
cost starting from $\mathsf T^n\iota_1(I)$ is a tail of that
nonnegative series and is finite.  Hence the total cost starting
from $R$ is finite.  On its support there is an $\varepsilon>0$
with $R\succeq\varepsilon P$.  Positivity makes total cost from
$P$ finite as well.

The positive cost floor then gives
$\sum_{n\ge0}\tr(\mathsf T^n(P))<\infty$.
Order domination by $P$ shows that the restricted powers tend to
zero on every element of the corner.  In finite dimension their
spectral radius is strictly less than one: an eigenvalue of modulus
at least one would contradict this convergence.  Thus the restricted
geometric series and its adjoint converge.  The finite positive
cost operator is
\[
 W=\sum_{n=0}^\infty (\mathsf T_P^*)^n(PKP),
 \qquad \mathsf T_P^*(Z)=P\mathsf T^*(Z)P.
\]
This proves the assertion on every reachable block, including those
whose exposure probability is small.
\end{proof}

\begin{theorem}[Exact finite support and Lyapunov certificate]
\label{thm:r79-exact-certificate}
Fix $M$, a controller $\gamma$, a model $\theta$, and $b<\infty$.
Its cost is at most $b$ for every initial density matrix if and only
if there exist block projections $P$ and Hermitian block operators
$W$ satisfying
\begin{align}
 &P_k=P_k^*=P_k^2,
 \qquad P_1=I,
 \qquad P\mathsf T^\theta_\gamma(P)P
                      =\mathsf T^\theta_\gamma(P),
                    \label{eq:r79-support-certificate}\\
 &W=PWP\succeq0,
 \qquad
 W\succeq PK^\theta_\gamma P
            +P(\mathsf T^\theta_\gamma)^*(W)P,
                    \label{eq:r79-lyapunov-certificate}\\
 &W_1\preccurlyeq bI.\label{eq:r79-initial-budget}
\end{align}
The projections may be chosen as the least reachable supports in
Lemma~\ref{lem:r79-reachable-corner}.  With that choice the minimal
positive cost witness satisfies equality in
\eqref{eq:r79-lyapunov-certificate}, and $\tr(W_1\rho)$ is the
actual expected cost from $\rho$.
\end{theorem}
\begin{proof}
Suppose first that the certificate holds.  Invariance keeps every
nonterminal state, starting from $\iota_1(\rho)$, inside its corner.
Pair \eqref{eq:r79-lyapunov-certificate} with such a state and
iterate.  For every $N$,
\[
 \sum_{n=0}^{N-1}
   \langle K^\theta_\gamma,
      (\mathsf T^\theta_\gamma)^n\iota_1(\rho)\rangle
 +\langle W,(\mathsf T^\theta_\gamma)^N\iota_1(\rho)\rangle
 \le\tr(W_1\rho)\le b.
\]
Drop the nonnegative remainder and use monotone convergence.  This
gives the asserted total-cost bound.  Its positive cost floor also
bounds the expected number of calls by $b/c_*$, so success occurs
almost surely.

Conversely, the assumed uniform initial bound makes the construction
in Lemma~\ref{lem:r79-reachable-corner} applicable.  Take its least
support $P$ and its convergent cost series $W$.  The series is
supported on $P$ and solves the Lyapunov equation there.  Pairing it
with any initial density matrix gives exactly
\eqref{eq:r79-total-cost}; hence $\tr(W_1\rho)\le b$ for every
$\rho$, which is equivalent to \eqref{eq:r79-initial-budget}.
Its minimality among positive supersolutions follows by iterating
the inequality from zero.  This establishes both directions without
postulating finite costs at inaccessible labels or subspaces.
\end{proof}

\subsection{Attainment and exact synthesis at every memory budget}

\begin{theorem}[Robust quantum finite-memory frontier]
\label{thm:r79-quantum-frontier}
Assume the equivalent feasibility conditions of
Theorem~\ref{thm:r79-dark-feasibility}.  Then for every $M\ge1$ the
value $V_M$ in \eqref{eq:r79-general-frontier} is finite and attained
by a controller in $\mathcal C_M$.  It is nonincreasing in $M$.
For every finite $b$, the following is a necessary and sufficient
condition for $V_M\le b$:
\begin{equation}\label{eq:r79-robust-budget-quantifiers}
 \exists\,\gamma\in\mathcal C_M\quad
 \forall\,\theta\in\Theta\quad
 \exists\,P^\theta,W^\theta\quad
 \text{\eqref{eq:r79-support-certificate}--
       \eqref{eq:r79-initial-budget} hold.}
\end{equation}
The controller variables are common to all models.  The certificate
variables may depend on the fixed model; they are proof witnesses
and are not supplied as information to the controller.

Suppose in addition that $\Theta\subseteq\mathbb R^r$ is a compact
semialgebraic set defined over the real algebraic numbers, and the
real and imaginary matrix coefficients of the primitive maps and
costs are polynomial or rational functions of $\theta$ with algebraic
coefficients and no poles on $\Theta$.  For each fixed $M$, feasibility
at any algebraic budget is decidable by finite real-algebraic
elimination.  The value $V_M$ is algebraic, and one attaining
controller can be chosen with algebraic command and update
probabilities and computed by a finite exact procedure.
\end{theorem}
\begin{proof}
The one-row policy of Theorem~\ref{thm:r79-dark-feasibility} can be
represented in $\mathcal C_M$, giving $V_M\le C_0$.
For fixed $\theta,\rho,N$, the truncated sum in
\eqref{eq:r79-total-cost} is a continuous function of the finitely
many controller probabilities.  The full robust objective $J$ is
the supremum of these continuous functions over $\theta,\rho,N$.
It is therefore lower semicontinuous on the compact space
$\mathcal C_M$.  Its finite minimum is attained.  This argument
does not require a uniform upper bound for the potential on every
label.  Padding by unused labels gives monotonicity.

The fixed-model certificate theorem applies to the same controller
at every $\theta$.  This proves both directions of
\eqref{eq:r79-robust-budget-quantifiers}; attainment identifies the
budget condition with $V_M\le b$, including $b=V_M$.

For the effective statement, express all complex matrices in real
and imaginary coordinates.  Projection identities, corner
invariance and controller normalization are polynomial identities
or inequalities.  Positive semidefiniteness is expressed by the
nonnegativity of all principal minors of each Hermitian block.
Rational coefficients introduce no difficulty, since their
denominators have no zeros on the specified model set and may be
handled by polynomial equations for their inverses.  Thus
\eqref{eq:r79-robust-budget-quantifiers} is a finite first-order
formula over the real algebraic numbers.  Its quantifier order is
exactly the displayed order, so elimination does not give the
controller access to $\theta$.

Elimination produces the feasible budget set.  It is the nonempty
ray $[V_M,\infty)$ by the budget equivalence and attainment.  A
finite endpoint of a one-dimensional semialgebraic set defined over
the real algebraic numbers is algebraic.  Substituting this endpoint
in the quantified formula leaves a nonempty semialgebraic set of
controller probabilities defined over the same field.  Exact
sampling supplies algebraic probabilities in that set.  This is
a finite exact procedure; no polynomial complexity, numerical tie
certification, or previously executed general elimination is claimed.
\end{proof}

For a finite model set, \eqref{eq:r79-robust-budget-quantifiers} is
simply a finite family of support and cost certificates with shared
controller variables.  For a compact infinite set its universal
quantifier is essential.  The theorem does not assert that a
least-favorable model attains the outer supremum, that $V_M$ has a
general closed formula, or that the unrestricted-history optimum
is reached at a finite label budget.  It also does not identify
finite algebraic synthesis with efficient computation.

With the single constant terminal decoder, event stopping keeps the
same total-label frontier $V_M$.  State stopping uses one disjoint
terminal label, giving the live-label frontier $V_M$ at $M+1$ total
labels.  These counts include the prescribed initial label.

\subsection{Why fixed models and reachable supports cannot be removed}

\begin{example}[Information remains after a physical reset]
\label{ex:r79-hidden-model-memory}
Let $d=1$, let there be a single nonterminal report, and let
$\Theta=\{0,1\}$.  Each call costs one.  The success probabilities
of commands $a,b$ are respectively
\[
 (1/2,1/10)\quad\hbox{under model }0,
 \qquad (1/10,1/2)\quad\hbox{under model }1.
\]
All success and failure probabilities are strictly positive, and
there is only one physical reset state.

A one-label controller uses one command row, choosing $a$ with
probability $p$.  Its two mean costs are
\[
 \frac1{1/10+(2/5)p},\qquad
 \frac1{1/2-(2/5)p}.
\]
Their maximum is minimized at $p=1/2$, with value $10/3$.
Now use three counted rows: an initial uniform row, a pure $a$ row,
and a pure $b$ row.  After the first failure choose the row for the
opposite command, and thereafter alternate the two pure rows.
For model $0$, starting with $a$ gives mean cost
$ (1+1/2)/(1-(1/2)(9/10))=30/11$; starting with $b$ gives
$38/11$.  Model $1$ exchanges these numbers.  Thus the prescribed
uniform initial row followed by counted alternation has robust cost
$34/11<10/3$, an exact improvement of $8/33$.

In contrast, a deterministic update depending only on the common
report must always select one of these same three rows.  Selecting
the uniform row gives cost $10/3$; selecting either pure row after
the first call gives worst-model cost $1+(7/10)10=8$.  Therefore
no such report-only update dominates this controller with its same
rows.  The palette theorem for one known reset instrument has not
failed: its physical resetting hypothesis does not erase information
about a fixed hidden model, and that theorem does not include this
uncertainty quantifier.
\end{example}

The same example has a known-model, genuinely nonreset quantum
realization.  On $\mathbb C^2$ use single failure Kraus operators
\[
 A_a=\operatorname{diag}(1/\sqrt2,3/\sqrt{10}),\qquad
 A_b=\operatorname{diag}(3/\sqrt{10},1/\sqrt2),
 \qquad G_a=G_b=I.
\]
Both failure operators are invertible, so their posterior states
depend on the input.  The success effects are
$I-A_a^*A_a$ and $I-A_b^*A_b$.  All expected-cost operators in the
displayed controllers are diagonal.  The two basis states therefore
give exactly the two scalar cost calculations above, and arbitrary
coherences do not change them.  In this known instrument the
all-initial-state one-label value is again $10/3$, while the
three-label construction attains $34/11$.  This supplies an explicit
full-rank nonreset instance of the larger certificate class.

\begin{example}[Finite initial cost with unbounded continuation witnesses]
\label{ex:r79-unbounded-reachable-witness}
Let $d=1$, $\Theta=[0,1]$, and keep one nonterminal report and unit
cost.  The failure probabilities of $a,b$ are
\[
 f_a(\theta)=\theta(1-\theta),\qquad f_b(\theta)=1-\theta.
\]
Consider the two-label controller which starts with $a$ and, after
failure, uses $b$ forever.  At $\theta>0$ its cost is
$1+\theta(1-\theta)/\theta=2-\theta$; at $\theta=0$ its first
call succeeds surely, so its cost is $1$.  Its robust cost is
therefore $2$, although no parameter attains that supremum.

For $0<\theta<1$, the second label is reachable and its true
continuation cost is $1/\theta$, which is unbounded as
$\theta\downarrow0$.  At $\theta=0$ that label is both unreachable
and dark.  A Lyapunov inequality on the entire two-label space would
then demand $W_2\ge1+W_2$, which is impossible.  The correct support
certificate instead has $P_1=1,P_2=0$ at this endpoint.  At $0<\theta<1$ it has both projections equal to one and the exact cost
witnesses $W_1=2-\theta,W_2=1/\theta$.  At $\theta=1$ the second label is again unreachable; either its zero support or the larger transient support gives a valid certificate.

Every primitive model in this example satisfies the feasibility
theorem: command $a$ succeeds with probability at least $3/4$.
The obstruction is therefore neither infeasibility of the model
family nor a zero cost.  It is the unjustified requirement that
every unused controller state possess a global finite potential,
or that all model-dependent continuation witnesses have a uniform
bound merely because the initial cost does.
\end{example}

The product evolution, positive cost operators, and Lyapunov
verification have classical antecedents in finite-state controller
synthesis and quantum program analysis.  The present statement keeps
the counted controller, the arbitrary initial state, the fixed-model
quantifier, and support changes in one necessary-and-sufficient
certificate.  Its role is to extend the finite resource frontier
beyond reported resetting while preserving the stronger scalar
normal form, closed cover solutions and rigidity results in their
previously proved reset class.  It does not silently extend the
report-only normal form or solve the separate squared-sign minimax
risk and optional-commitment problems.
